\documentclass[12pt]{article}
\usepackage[english]{babel}
\usepackage[english]{isodate}
\usepackage[table,svgnames]{xcolor}
\usepackage{url}
\usepackage[utf8x]{inputenc}
\usepackage[T1]{fontenc}
\usepackage{longtable}
\usepackage{amsmath}
\usepackage{graphicx}
\usepackage{parskip}
\usepackage{fancyhdr}
\usepackage{vmargin}
\usepackage{hyperref}
\usepackage{pgfgantt}
\usepackage{xparse}
\usepackage{float}
\usepackage{tabularx}
\usepackage{titling}
\usepackage[nogroupskip,toc,acronym]{glossaries}
\usepackage{fancyhdr}
\usepackage[authoryear]{natbib}
\usepackage{enumitem}
%\bibliographystyle{fontysIEEtranN}

\pagestyle{fancy}
\fancyhf{}
\lhead{\includegraphics[width=0.2\linewidth]{TreewatchLogo.pdf}}
\rhead{Personal dossier Max van der Linden}
\rfoot{Page \thepage}

% alternate row colors for all tables
\definecolor{lightGrey}{rgb}{0.95,0.95,0.95}

\let\oldtable\table
\let\endoldtable\endtable
\renewenvironment{table}{\rowcolors{2}{lightGrey}{}\oldtable}{\endoldtable}

\let\oldtabular\tabular
\let\endoldtabular\endtabular
\renewenvironment{tabular}{\rowcolors{2}{lightGrey}{}\oldtabular}{\endoldtabular}

\let\oldtabularx\tabularx
\let\endoldtabularx\endtabularx
\renewenvironment{tabularx}{\rowcolors{2}{lightGrey}{}\oldtabularx}{\endoldtabularx}

\let\oldlongtable\longtable
\let\endoldlongtable\endlongtable
\renewenvironment{longtable}{\rowcolors{2}{lightGrey}{}\oldlongtable} {\endoldlongtable}

\graphicspath{{../img/}}

\renewcommand{\arraystretch}{1.5}

\DeclareDocumentCommand{\newdualentry}{ O{} O{} m m m m } {
  \newglossaryentry{gls-#3}{name={#5},text={#5\glsadd{#3}},
    description={#6},#1
  }
  \makeglossaries
  \newacronym[see={[Glossary:]{gls-#3}},#2]{#3}{#4}{#5\glsadd{gls-#3}}
}


\newlist{SMART}{description}{2}
\setlist[SMART]{leftmargin=12em,style=nextline}

\newlist{STAR}{description}{2}
\setlist[STAR]{leftmargin=8em,style=nextline}


%\loadglsentries{../glossary.tex}

\setmarginsrb{3 cm}{2.5 cm}{3 cm}{2.5 cm}{1 cm}{1.5 cm}{1 cm}{1.5 cm}

%\makeglossaries
\begin{document}

%%%%%%%%%%%%%%%%%%%%%%%%%%%%%%%%%%%%%%%%%%%%%%%%%%%%%%%%%%%%%%%%%%%%%%%%%%%%%%%%%%%%%%%%% Preface of the report

    \pagenumbering{roman} % Roman numerals for page counter


    % Title Page
    \begin{titlingpage}
        \begin{center}
            \begin{minipage}{\linewidth}
            \centering
            %University logo
            \includegraphics[width=0.3\linewidth]{FontysLogo.pdf}
            \par
            \vspace{3cm}
            %Thesis title
            {\uppercase
                {\Large Personal dossier\\ Max van der Linden \\ 2015 Sofa GTL \\ TreeWatch Project
            \par
            \vspace{3cm}}}
            \includegraphics[width=0.5\linewidth]{TreewatchLogo.pdf}
            \par
            \vspace{2cm}
            %Author's name
            {Max van der Linden\\
	         \par}
            \vspace{2cm}

            %Date
            \today
            \end{minipage}
        \end{center}
    \end{titlingpage}
    \clearpage

    \section*{Document information}
\addcontentsline{toc}{section}{Document information}
\renewenvironment{tabular}{\oldtabular}{\endoldtabular}
	\begin{tabular}{ll}
		\textbf{Document name:} & Personal dossier Max van der Linden\\
		\textbf{Document owner:} & Max van der Linden \\
		\textbf{Company/Organisation:} & Fleuren Baarlo \\
		\textbf{Contact person:} & Max van der Linden \\
		\textbf{Date:} & \today \\
		\textbf{Place:} & Fontys University of Applied Science Venlo \\
		\textbf{Author:} & \parbox[t]{5cm}{
		Max van der Linden\\ max.vanderlinden@student.fontys.nl\\ 2209349 \\}
	\end{tabular}
\renewenvironment{tabular}{\rowcolors{2}{lightGrey}{}\oldtabular}{\endoldtabular}

    %\printglossary[type=\acronymtype]
    %\printglossary
    \pagebreak

    %\listoffigures
    %\addcontentsline{toc}{section}{\listfigurename}
    %\listoftables
    %\addcontentsline{toc}{section}{\listtablename}
    %\pagebreak

    \tableofcontents
    \clearpage

%%%%%%%%%%%%%%%%%%%%%%%%%%%%%%%%%%%%%%%%%%%%%%%%%%%%%%%%%%%%%%%%%%%%%%%%%%%%%%%%%%%%%%%%% Main Part of the Report

    \pagenumbering{arabic}
    
    \section{Project management skills}
    
    This learning goal was created because at the beginning of the project I was elected as Project Manager and in my opinion it is the task of the project manager to ensure that the project runs smoothly. \\
    I choose this one, because also in the future I might end up as a project manager and this project/learning goal would make a good trial run.
    
	\subsection{SMART learning goal}
	\begin{SMART}
	    \item[Specific] As a project manager in the GreenTechLab Software Factory, my goal is to assure the project will run smoothly. (deadlines are met, deliverables are delivered, documentation is written and solve problems within the group)
	    \item[Measurable] Since all of the factors delay the project a good measure would be to see if all the deliverables are completed on time. It would also be good to receive feedback from fellow teammates to find out how they think about my performance. Using a questionnaire half way and at the end.
	    \item[Attainable] It should be attainable since the group already knows each other and therefor could work together very well.
	    \item[Relevant] If I ever want to be a project manager I will have some experience.
	    \item[Time-limited] By January 2016, this project will be cosed, therefor I will be able to manage the project until then.
	    \item[My complete goal] To improve my project management by January 2016.
	\end{SMART}
	
	\subsection{STAR Analysis}
	\begin{STAR}
	    \item[Situation] I was elected to be the project manager and therefore responsible for the project.
	    \item[Task] I should make sure everybody sticks to the deadlines and everyone follows the rules set at the start of the project.
	    \item[Action] At the beginning of the project I started by creating a document which stated the rules and roles. I then asked all team members if they agreed with these rules. Everyone did so now I had a handle to refer to if someone did not folow the rules. I also kept presence in the minutes to make sure that everyone was present and on time to make sure no time was lost because of people being absent.
	    \item[Result] 
	    You can have a look at the rules and roles document. Also the minutes show the presence of the team members during the weekly project meetings.
	    \item [Reflection]
	    The project was running relatively smoothly until the last week. Then we found out we were creating a final report in stead of a project dossier. This was a problem which will be further discussed in the problem solving learning goal. When not looking at the last week the project was doing very well. Everybody worked on their tasks properly, everyone kept to their deadlines and everyone was present and on time (some exceptions).
	\end{STAR}
	
	\section{Communication skills}
	This learning goal was created because I personally have problems with picking up the phone and calling people. This learning goal is about trying to remove this "fear" by facing it.
	
	\subsection{SMART learning goal}
	\begin{SMART}
	    \item[Specific] Get better at communicating with a customer, not get nervous anymore when contacting people over the phone.
	    \item[Measurable] Call at least 3 people during the SoFa.
	    \item[Attainable] There are a lot of people that need to work together to complete the SoFa, not just the group but also GreenTech Lab, VAA and Fleuren.
	    \item[Relevant] In my further carrier it would be helpful to be able to call people without hesitation.
	    \item[Time-limited] January 2016
	    \item[My complete goal] To improve my communication skills January 2016.
	\end{SMART}
	
	\subsection{STAR Analysis}
	\begin{STAR}
	    \item[Situation] In the project I was elected as the main contact point of the customer and therefore would need to make phone calls
	    \item[Task] Make sure that I have at least 3 times phone contact with people and ideally get over the phone "fear" 
	    \item[Action] During the project I did call the customer multiple times and also called other people. For instance I had phone contact with Han Fleuren about a possible conflict of interest between VAA ict and Fleuren. And also some phone calls about setting up meetings and about requirements with Yannick.
	    \item[Result]
	    No document result, however the making of phone calls is much easier for me.
	    \item[Reflection]
	    During the runtime of the project many people have been called by me and i have been called by many people (exceeding 10 phonecalls). In the beginning I noticed it was hard to pick up the phone and call for instance Yanning but towards the end of the project it got much easier for me to just pick up the phone and call any of the stakeholders in the project since I now had more experience with phone calls.
	\end{STAR}
	
	\section{Problem solving skills}
	This learning goal was created because as a project manager it is important to be able to solve problems within the group, and also outside the group. I chose this one because I think it might help me in my future job to be able to tackle problems in a project. The Sofa project is a perfect opportunity to test and improve on those skills.
	
	\subsection{SMART Learning Goal 3}
	\begin{SMART}
	    \item[Specific] Get better at solving problems in groups and/or projects.
	    \item[Measurable] If problems occur, after SoFa is over ask the involved parties how they feel the problem was handled using a questionnaire.
	    \item[Attainable]  Problems always occur. When they occur I should step in.
	    \item[Relevant] If I would work as a project manager I would run into problems. Therefore the SoFa would be a good practice.
	    \item[Time-limited] January 2016
	    \item[My complete goal] To improve my problem solving skills withing a group by January 2016.
	\end{SMART}
	
	\subsection{STAR Analysis 3}
	\begin{STAR}
	    \item[Situation] During the project I was elected as project manager and therefor responsible to solve problems that occur during the project.
	    \item[Task] Solve the problems that occur during the project.
	    \item[Action] In general the project did run pretty smoothly. However relatively big problems occured during the project.\\
	    The first problem was that after the meeting we had with van den Boorne Aardappelen, we had the feeling that the system that we were supposed to create already existed. If this was the case it would mean that the entire SoFa project would be useless and that was also the feeling that was in the group. To solve the problem I called Han Fleuren (the boss of Boomkwekerij Fleuren) since Yannick, our contact person, was on holiday. I explained to Han that we had the feeling that VAA ICT wanted to use us to replicate a system that already existed and that we had the feeling that they wanted to use us just for cheap labor so they could afterwards sell the system. Over the phone we then decided that we should have a meeting with Randy (consultant for VAA), Han and the whole SoFa group. I than had to contact Randy to invite him to the meeting without "attacking" him personally. Randy accepted the invite to the meeting and during the meeting we discussed the issue and solved it. The final outcome was more or less that the system should have more functionality than Crop-r and that we should view Crop-r as an example of what we would develop.
	    
	    The Second problem arrived towards the end of the project when we were almost finished with the creation of the final report. We were almost done so we decided to have a talk with an other sofa group and ask how they were doing and what they are working on. We then found out that we we had to create a project dossier. Since we had only one day left this was going to be a big problem which I had to solve. I first went to mister van Odenhoven to discuss the issue and he told us that we would anyway have to deliver a final dossier so that we would have to solve the problem somehow. So I then started out by calling mister Jacobs, he didn't answer so I sent a mail explaining the problem. Later that day mister Jacobs called back to discuss the problem. I asked him what we could do since without a personal dossier we would probably fail the module. I told him we could do it if we had more time but that he would have to verify if we could get the final deadline of the project dossier moved. Mister Jacobs then proceeded to call mister van Odenhoven and he gave us a new deadline. In the mean time we already started by defining which products needed to be in the dossier. When the week after that to finish the project dossier and personal dossier. Still all within the extended deadline period.
	    \item[Result] The two problems which were solved. Basically the working prototype of the app and the project dossier are deliverables that show the problem solving.
	    \item[Reflection] I realy learned a lot from the problems that occured. It thought me to stay calm and just discuss the problems that occur. Also if you can't solve a problem on your own ask for help.
	\end{STAR}
	\clearpage
	
	\section{Development}
	Upfront I didn't register development as a learning goal since development would not be a big part of my future job. Nevertheless I did help with the implementation and also did some implementation myself and in pair programming.
	
	\subsection{Geohelper}
	In order to provide functionality to the app I created a so called Geohelper (class can be found in the geohelper subsection of the important code document). The class has multiple different methods. The method CalculateBoundingBox for instance calculates the width and height of a field and also provides a center which than can be used to position the view of the app so that when a field is clicked it shows the whole field in one view nice and centered. In order to do this the distance (width and height) needs to  be measured in meters. This is done with the DistanceToMeters method.
	
	\subsection{Geofencing}
	During the implementation of the Geofencing functionality I worked together with Martijn in a pair. 
	
	\subsection{Mock data}
	Originally my task was to provide the app with mock data. This meant to get coordinates into the app which then would draw certain overlays, on the map, on specific locations. The mock data enabled the whole group to easily see the main functionality of the app (visualizing overlays on the map). Eventually these coordinates were replaced with real data we received from Yannick so that we could have the actual fields represented on the map.
	
	
	%\clearpage
	
	%\bibliography{biblist}
	
	\end{document}
